\section{Conclusion}
\label{sec:conclusion}

We have argued, and shown on frozen EEG artifacts, that a decodable subject/domain signal is not by
itself evidence of a harmful functional shortcut. Measurable leakage magnitude does not certify
functional reliance --- task-head alignment is closer to reliance, though not a validated estimator.
Subject signal is erasable, but erasure strength does not certify a target benefit: no eraser achieves
a proven gain across 40 held-out-target cells, and aggressive erasure can harm the target. Branch load
matters --- a fusion backbone's spatial branch is load-bearing --- but branch-local leakage and
reliance are unmeasured, and the instrument to measure them does not currently exist.

\begin{center}
\emph{Measurable is not reliance. Erasable is not beneficial. Branch load matters, but branch-local
shortcut claims require missing instruments.}
\end{center}

The contribution is an audit framework, not a new domain-generalization method: the FSR ladder makes
explicit what measurement is required to support what control conclusion, and the two boundary findings
plus one honestly-blocked question show current observables do not certify shortcut reliance.
Tables~\ref{tab:claims}--\ref{tab:routes} record every claim's status and every route's inclusion, so
the boundary between what the evidence supports and what it does not is auditable.

\begin{table}[t]
\centering\small
\begin{tabular}{llp{7.6cm}}
\toprule
ID & Status & Claim \\
\midrule
C1 & READY\_WITH\_CAVEAT & Measured leakage magnitude does not certify reliance. \\
C2 & READY\_WITH\_CAVEAT & Task-head alignment is closer to reliance than raw leakage in frozen CIGL R3. \\
C3 & READY & Subject signal is erasable. \\
C4 & READY & Erasure strength does not certify target benefit. \\
C5 & READY\_WITH\_CAVEAT & Random-k falsifies nonspecific NLL movement. \\
C6 & READY & Spatial branch is load-bearing. \\
C7 & READY & Branch-local leakage/reliance is missing. \\
C8 & READY & CMI-control remains closed. \\
C9 & SUPPORT\_ONLY & TTA-Control is positive but non-CMI. \\
C10 & READY & FSR is an audit framework, not a new DG method. \\
\bottomrule
\end{tabular}

\caption{Claim ledger. Statuses are derived from the frozen result artifacts, not hand-set.}
\label{tab:claims}
\end{table}

\begin{table}[t]
\centering\small
\begin{tabular}{llll}
\toprule
Route & Predictor & Endpoint & Inclusion \\
\midrule
CIGL & L1|L2|L4 & L5|L6 & INCLUDED \\
TOS\_mean\_scatter & L1|L3|L4 & L6 & INCLUDED \\
TOS\_LEACE & L1|L3|L4 & L6 & INCLUDED \\
TOS\_INLP & L1|L3|L4 & L6 & INCLUDED \\
TOS\_RLACE & L1|L3|L4 & L6 & INCLUDED \\
TOS\_random\_k & L3 & L6 & INCLUDED \\
\midrule
\multicolumn{4}{l}{31 other routes: SUPPORT/BOUNDARY/PROTOCOL/BACKGROUND\_ONLY (no target-label fit enters any RQ)} \\
\bottomrule
\end{tabular}

\caption{Route inclusion. Only routes with a predictor$+$endpoint pairing enter quantitative tests; no
target-label-fit route does.}
\label{tab:routes}
\end{table}
